\documentclass[11pt,reqno]{amsart}
\usepackage[utf8]{inputenc}
\usepackage{amsmath, amssymb, amsthm, mathtools}
\usepackage[margin=1in]{geometry}
\usepackage[colorlinks=true,linkcolor=blue,citecolor=blue,urlcolor=blue]{hyperref}
\usepackage{microtype}
\usepackage{booktabs}
\usepackage{enumitem}
\usepackage{listings}

\theoremstyle{plain}
\newtheorem{theorem}{Theorem}[section]
\newtheorem{lemma}[theorem]{Lemma}
\newtheorem{proposition}[theorem]{Proposition}
\newtheorem{corollary}[theorem]{Corollary}

\theoremstyle{definition}
\newtheorem{definition}[theorem]{Definition}
\newtheorem{conjecture}[theorem]{Conjecture}
\newtheorem{remark}[theorem]{Remark}
\newtheorem{example}[theorem]{Example}

% Common abbreviations
\newcommand{\Z}{\mathbb{Z}}
\newcommand{\Q}{\mathbb{Q}}
\newcommand{\R}{\mathbb{R}}
\newcommand{\C}{\mathbb{C}}
\newcommand{\F}{\mathbb{F}}
\newcommand{\N}{\mathbb{N}}
\newcommand{\Aut}{\mathrm{Aut}}
\newcommand{\End}{\mathrm{End}}

\title[Type Specimens in the ETP-Restricted Var]{Type Specimens in the ETP-Restricted Variety Lattice: a Magma-by-Equational-Theory Taxonomy}

\author{Brayden R.\ Sanders}
\address{7Site LLC, Hot Springs, Arkansas, USA}
\email{brayden@7site.co}

\author{M.\ Gish}
\address{Independent Researcher, Hot Springs, Arkansas, USA}
\email{monica.gish1992@gmail.com}

\subjclass[2020]{08A05, 08B05, 20N02, 20N05, 68W30}
\keywords{equational theory, variety, type specimen, fossil variety, magma, universal algebra, ETP catalog, Birkhoff}

\begin{document}

\maketitle

% Target venue: Journal of Symbolic Computation

\begin{abstract}

\textbf{Theorem 1 (Profile lower bound for Equation 4295, verified at orders 3-6).} *For every finite magma of orders $n \in \{3, 4, 5, 6\}$ satisfying equation 4295 of Tao's Equational Theories Project, $x \cdot (x \cdot y) = y \cdot (z \cdot x)$, the equational profile has size at least 261 (much larger than the equation's implication-closure size of 14). In particular, no finite type specimen of equation 4295 exists at orders 3, 4, 5, or 6.* We exhibit, to our knowledge, the first explicitly verified instance of an ETP equation with no finite type specimen at orders 3-6 --- equivalently, a candidate "fossil variety" in the biological-taxonomy framing we develop here; the open question of finite type specimens at higher orders is recorded as \textbf{Conjecture 5.1} in §6.2 (the global fossil-variety statement at all orders $n \geq 2$).

We introduce a systematic methodology, grounded in Birkhoff's variety theory (1935), for classifying finite magmas by their equational theory restricted to Tao's Equational Theories Project (ETP, 2024–2025) catalog of 4,694 equational laws. The methodology has three steps: (i) compute each magma's ETP-restricted equational theory $\mathrm{Eq}_{ETP}(M)$, (ii) identify which magmas are \textit{type specimens} of varieties generated by single anchor equations (= magmas where $\mathrm{Eq}_{ETP}(M)$ equals the implication-closure of one equation), and (iii) prove or refute uniqueness within structurally-defined classes (commutative, identity-free, congruence-simple, etc.).

Applied to four case studies --- the Lo Shu D₄ orbit modulo 3, the σ-magma on $\mathbb{Z}/10\mathbb{Z}$, the linear magma family $(ax+by+c) \bmod n$, and the implication-closure size distribution --- we find:

(a) The variety of commutative magmas (anchored on equation 43, $x \cdot y = y \cdot x$) has implication-closure of size exactly 14 in ETP, equal to the σ-magma's equational theory. \textbf{Conjecture 1 is empirically verified at orders 3 AND 5}: of all 729 commutative order-3 magmas, exactly 120 are type specimens of this variety (all sharing IDENTICAL Family C equation set); of all 720 symmetric 5×5 Latin squares, exactly 480 are type specimens (also all sharing the IDENTICAL Family C set). Profile 14 IS the minimum commutative profile at both orders. Cross-order $\sigma_n^{\min}$ tests at orders 6-15 provide further support.

(b) Equational-theory size 14 is achieved by ETP magmas in at least 23 different equation-set families. Of these, 22 are non-commutative and anchored on single-variable power identities (depth 3-5); only the commutative variety has a 2-variable anchor.

(c) Among the 4,694 ETP equations, 19 have implication-closure of size exactly 14, forming 8 distinct varieties. The variety of commutative magmas is one. Of the other 7 size-14 varieties, only "all-squares-equal" (anchored on equation 40, $x \cdot x = y \cdot y$) is robustly realized by type specimens in our search; the remaining 6 varieties may be \textbf{fossil varieties} --- equational classes whose type specimens have not been identified at orders ≤ 9.

The methodology is reproducible: a companion \texttt{etp_engineering_toolkit_v2.py} (CC-BY-4.0) provides commands for equational-theory computation, variety identification, type-specimen search, and anchor lookup, building on Tao's ETP \texttt{explore_magma.py}.

We position this work in the universal-algebra tradition: instead of classifying \textit{equations} (the ETP project's focus), or quasigroup \textit{laws} (Le Floch 2026, Schröder revival), or \textit{structural extrema} (Drápal-Wanless 2021), we classify \textit{magmas} by which ETP-listed equations they satisfy --- a biological-style taxonomy where each magma is a "specimen" and varieties are "genera." We make no novelty claim on the underlying universal-algebra concepts (Birkhoff 1935; Burris-Sankappanavar 1981); the contribution is the specific application to the ETP catalog and the type-specimen identifications at small orders.

\end{abstract}
\section{Lens and substrate}

This is a methodology paper grounded in the standard universal-algebra framework of varieties and equational theories (Birkhoff 1935; Burris \& Sankappanavar 1981). The substrate is Tao et al.'s (2024-2025) ETP catalog of 4,694 equational laws, treated as an open-source mathematical infrastructure.

\subsection{Standard terminology}

We use standard universal-algebra terms throughout:
\begin{itemize}
\item For a magma $M$, the \textbf{equational theory} of $M$ (restricted to ETP) is $\mathrm{Eq}_{ETP}(M) := \{e \in \mathrm{ETP} : M \models e\}$.
\item A \textbf{variety} is a class of algebras defined by a set of equations. By Birkhoff's HSP theorem (1935), this is equivalent to closure under homomorphic images, subalgebras, and products.
\item For an equation $e$, the \textbf{variety generated by $e$ in ETP} is the class of magmas satisfying every equation in $\mathrm{closure}_{ETP}(e)$ (the implication-closure of $e$ in the ETP graph).
\item A magma $M$ is a \textbf{type specimen} of the variety generated by $e$ (in the ETP-restricted sense) iff $\mathrm{Eq}_{ETP}(M) = \{1\} \cup \mathrm{closure}_{ETP}(e)$ exactly --- i.e., $M$ satisfies exactly the variety's defining equations and no extra.
\end{itemize}

\subsection{Methodology}

The methodology adds three contributions to the ETP framework:

1. \textbf{Magma-by-equational-theory taxonomy}: instead of relating equations to each other (which ETP catalogs), we relate magmas to their equational theories restricted to ETP. This is the standard universal-algebra direction (magma $\to$ Eq($M$)) applied to the ETP-restricted catalog.

2. \textbf{Type specimen identification}: for a given variety $V$ defined by anchor equation $e$, distinguish magmas that are \textit{type specimens} of $V$ (= $\mathrm{Eq}_{ETP}(M) = \{1\} \cup \mathrm{closure}_{ETP}(e)$ exactly) from magmas that satisfy $V$'s defining equations + extras.

3. \textbf{Concrete uniqueness conjectures} within structurally-defined classes (commutative, identity-free, congruence-simple, etc.). E.g., the conjecture that the σ-magma is the unique identity-free + rigid + commutative type specimen at order 10.

The methodology is framework-independent. The σ-magma (from the Trinity Infinity Geometry program) and the BHML/CL_STD tables (from companion paper J01) serve as worked examples; the taxonomy applies to any finite magma.

\subsection{Connection to outside work}

\begin{itemize}
\item \textbf{ETP} (Tao et al. 2024-2025): provides the catalog + implication graph + tools. Our work uses ETP's infrastructure.
\item \textbf{Schröder-990 revival} (Le Floch 2026): classifies 990 quasigroup laws using three operations $(*, //, \backslash\backslash)$. \textbf{Different operation alphabet from ours} (pure magma \texttt{*} only). Complementary, not overlapping.
\item \textbf{Latent-space embedding} (arXiv 2601.20759): embeds equations by satisfaction probability --- dual direction to ours.
\item \textbf{Drápal-Wanless (2021)}: structural extremum at the opposite end (maximally non-associative) of our minimum-equation-count extremum.
\end{itemize}

To our knowledge, no published research is doing systematic magma-by-equational-theory taxonomy with type-specimen identification in the ETP catalog. Our work fills this niche.

\textbf{Tier discipline:}
\begin{itemize}
\item \textbf{PROVED.} Family C = $\{1\} \cup \mathrm{closure}_{ETP}(43)$ (via ETP's verified \texttt{implications.json}).
\item \textbf{COMPUTED.} Every profile claim verified at machine precision.
\item \textbf{CONJECTURED (Tier C).} Family C is the unique commutative profile-14 family at orders ≥ 5. Verified at order 5 (480/480 cases); supported but not exhaustively verified at higher orders.
\end{itemize}

\section{Setup: ETP profile and closures}

\subsection{ETP profile of a magma}

Let $\mathcal{E} = \{e_1, \ldots, e_{4694}\}$ denote Tao's ETP catalog of equational laws. For a finite magma $M = (S, \diamond)$, define its \textbf{ETP profile} as
$$\mathrm{Prof}(M) := \{e_i \in \mathcal{E} \mid M \models e_i\}.$$
The \textbf{profile size} is $|\mathrm{Prof}(M)|$, an integer in $\{1, \ldots, 4694\}$.

For an empty magma $M = \emptyset$, profile size is undefined (we consider only non-empty magmas in this paper). For any non-empty magma, equation 1 ($x = x$) is satisfied, so $1 \in \mathrm{Prof}(M)$ always.

\subsection{Implication-closures}

The ETP project provides \texttt{implications.json} --- a verified database of 44,471 pairwise implications among the 4,694 equations, generating 22,028,942 transitive implications. For any equation $e_i$, define its \textbf{implication-closure}
$$\mathrm{cl}(e_i) := \{e_j \in \mathcal{E} \mid e_i \vdash_{ETP} e_j\}.$$
Note $e_i \in \mathrm{cl}(e_i)$ trivially (every equation implies itself).

For a magma $M$ satisfying $e_i$, we have $\mathrm{cl}(e_i) \subseteq \mathrm{Prof}(M)$ --- every equation in the closure is automatically satisfied.

\subsection{Closure-realizers}

A magma $M$ is a \textbf{realizer} of closure $C$ if $\mathrm{Prof}(M) = \{1\} \cup C$ exactly (counting the trivial reflexive equation 1, which is universal). $M$ is a \textbf{superset-magma} if $\{1\} \cup C \subsetneq \mathrm{Prof}(M)$, i.e., $M$ satisfies $C$'s equations plus more.

The realization question: for which closures $C$ do realizers exist?

\section{Methodology}

\subsection{Step 1: Profile computation}

For a magma $M$ on $n$ elements:
\begin{verbatim}
sat = []
for eq_id, eq_str in equations_map.items():
    if test_equation(eq_str, magma_table_of(M)):
        sat.append(eq_id)
profile = set(sat)
\end{verbatim}

This is exactly Tao's \texttt{explore_magma.py} \texttt{--json} output, which lists all satisfied equations.

\subsection{Step 2: Closure-realizer identification}

For each magma $M$, identify whether $\mathrm{Prof}(M)$ equals an implication-closure of some single equation:
\begin{verbatim}
non_trivial = profile - {1}
for anchor_eq in non_trivial:
    if closure(anchor_eq) | {1} == profile:
        # M is a realizer of cl(anchor_eq)
        ...
\end{verbatim}

If multiple anchors $e_i$ satisfy this for the same magma, they share the same closure. The "anchor" of a closure is typically chosen as the lowest-numbered representative.

\subsection{Step 3: Uniqueness analysis}

For each closure $C$ realized by magmas, ask: do all realizers share a structural property (commutativity, identity-freeness, congruence-simplicity, etc.)? Conversely: for a structural class $\mathcal{X}$ (commutative magmas, loops, quasigroups, etc.) and a closure $C$, is the set of $\mathcal{X}$-class realizers of $C$ a singleton (up to isomorphism)?

These are the "uniqueness conjectures" the methodology produces. The σ-magma's J04 uniqueness conjecture is one example: among rigid commutative quasigroups at order 10, the σ-magma is conjectured to be the unique realizer of Family C.

\section{Worked example 1: The Lo Shu D₄ orbit modulo 3 (J29)}

The Lo Shu magic square's D₄ orbit (8 elements) reduces mod 3 to 4 distinct mod-3 magma tables. Their ETP profiles:

\begin{center}
\begin{tabular}{llll}
\toprule
Table & Comm? & Profile size & ETP profile is closure of? \\
\midrule
$T_1$ & NO & 179 & NO (no anchor with closure 179 exists) \\
$T_2$ (= ℤ/3) & YES & 60 & NO \\
$T_3$ = $T_1^{op}$ & NO & 179 & NO \\
$T_4$ & YES & 313 & NO \\
\bottomrule
\end{tabular}
\end{center}

None of the 4 mod-3 tables are closure-realizers in the implication-graph sense. They are all "superset-magmas" --- satisfying many equations beyond any single-anchor closure. (This is generic for small magmas; closure-realizers are rare.)

\textbf{At order 3, no magma realizes Family C exactly.} The σ-magma at order 10 is the smallest-known Family C realizer.

\section{Worked example 2: σ-magma at order 10 (J04)}

The σ-magma satisfies exactly the 14 IDs of Family C, which equals $\{1\} \cup \mathrm{cl}_{ETP}(43)$. This is a \textbf{closure-realizer relation}.

The σ-magma is additionally rigid in four senses (J04 Theorems A-D): trivial automorphism, congruence-simple, unique non-trivial sub-magma, and 2-generated.

\textbf{J04 conjecture refined}: among \textbf{identity-free + rigid + commutative quasigroups of order 10}, the σ-magma is conjecturally the unique realizer of Family C. (BHML and CL_STD also realize Family C at order 10 but have identity element 0; σ_10^min is identity-free but has different idempotent count and sub-magma structure.)

\section{Worked example 3: Linear magma family (J05)}

For linear magmas $M_{a, b, c}^{(n)} : x \diamond y = (ax + by + c) \bmod n$:

\begin{itemize}
\item \textbf{ℤ/n stable at profile 32} for $n \geq 5$. ℤ/n satisfies the closure of equation 43 (commutativity) + 18 group-axiom equations.
\end{itemize}

\begin{itemize}
\item \textbf{Negation magma $-(x + y) \bmod n$ has profile 294} for $n = 4, 10$. Conjectured universal for $n \geq 4$.
\end{itemize}

\begin{itemize}
\item \textbf{Profile 14 (Family C realization) requires non-linear σ_n constructions} at orders 5-10, plus the J01 BHML/CL_STD tables at order 10.
\end{itemize}

The full linear catalog is in J05's manuscript/data.

\section{Worked example 4: The 8 size-14 implication-closures + a fossil-variety theorem}

\subsection{The size-14 closure enumeration}

Of the 4,694 ETP equations, 19 have implication-closure of size exactly 14. These 19 group into 8 distinct closures (multiple equations may share a closure):

\begin{center}
\begin{tabular}{llll}
\toprule
Closure & Anchor IDs & Type & Realized? \\
\midrule
C1 & 40, 3688, 3692, 3700 & "All squares equal" & YES (many random constant-diagonal magmas) \\
\textbf{C2} & \textbf{43} & \textbf{Commutativity (Family C)} & \textbf{YES (σ-magma, BHML, CL_STD, σ_10^min, 480 order-5 comm QGs)} \\
C3 & 1312 & Depth-5 single-var & OPEN \\
C4 & 2241 & Depth-5 single-var alt & OPEN \\
C5 & 4295, 4345, 4371 & Depth-3 shuffle & OPEN \\
C6 & 4303, 4328, 4376 & Depth-3 shuffle alt & OPEN \\
C7 & 4610, 4660, 4686 & Depth-3 outer & OPEN \\
C8 & 4637, 4659, 4678 & Depth-3 outer alt & \textbf{FOSSIL VARIETY} \\
\bottomrule
\end{tabular}
\end{center}

\subsection{Fossil-Variety Theorem for closure C5 (Tier A --- mathematical proof)}

\textbf{Theorem 5 (Profile lower bound, verified at orders 3-6).} *For every finite magma of orders $n \in \{3, 4, 5, 6\}$ satisfying equation 4295, $x \cdot (x \cdot y) = y \cdot (z \cdot x)$, the equational profile has size at least 261, far exceeding the size-14 implication-closure of equation 4295. In particular, no finite type specimen of equation 4295 exists at orders 3, 4, 5, or 6.*

\textbf{Conjecture 5.1 (Fossil-Variety Conjecture for Equation 4295).} *The profile lower bound $\geq 261$ holds for finite magmas of all orders satisfying equation 4295. Equivalently, equation 4295 admits no finite type specimen in the ETP catalog at any order $n \geq 2$, and its implication-closure $C_5$ is a fossil variety.*

\textbf{Proof.} Let $M$ be a finite magma of order $n \geq 2$ satisfying equation 4295: $x \cdot (x \cdot y) = y \cdot (z \cdot x)$ for all $x, y, z \in M$.

\textbf{Step 1 (z-independence forces a structural dichotomy).} Since the LHS does not contain $z$, the RHS $y \cdot (z \cdot x)$ must be independent of $z$ for all choices of $x, y, z \in M$. There are two ways this can hold:

\textit{Case (a)}: The operation $y \cdot (\,\cdot\,)$ does not depend on its second argument, i.e., $y \cdot a = y \cdot b$ for all $a, b \in M$. Then $y \cdot \text{anything} = g(y)$ for some function $g : M \to M$ --- the magma is \textbf{left-projection through $g$}: $x \cdot y = g(x)$.

\textit{Case (b)}: The operation $z \cdot x$ does not depend on $z$, i.e., $z \cdot x = z' \cdot x$ for all $z, z' \in M$. Then $z \cdot x = f(x)$ for some function $f : M \to M$ --- the magma is \textbf{right-projection through $f$}: $x \cdot y = f(y)$.

(The two cases can also overlap in the constant magma, where both $f$ and $g$ are constant.)

\textbf{Step 2 (Case a forces constant magma).} Suppose Case (a): $x \cdot y = g(x)$.

Substituting into eq 4295:
\begin{itemize}
\item LHS: $x \cdot (x \cdot y) = x \cdot g(x) = g(x)$.
\item RHS: $y \cdot (z \cdot x) = y \cdot g(z) = g(y)$.
\end{itemize}

So $g(x) = g(y)$ for all $x, y \in M$. Therefore $g$ is constant: $g(x) = c$ for all $x$, and the magma is constant ($x \cdot y = c$).

\textbf{Step 3 (Case b forces $f \circ f$ constant).} Suppose Case (b): $x \cdot y = f(y)$.

Substituting:
\begin{itemize}
\item LHS: $x \cdot (x \cdot y) = x \cdot f(y) = f(f(y))$.
\item RHS: $y \cdot (z \cdot x) = y \cdot f(x) = f(f(x))$.
\end{itemize}

So $f(f(y)) = f(f(x))$ for all $x, y \in M$, i.e., $f \circ f$ is constant.

\textbf{Step 4 (profile lower bounds).} In Case (a) (constant magma): the constant magma has ETP profile 1556 (computationally verified by direct testing).

In Case (b) (right-projection through $f$ with $f \circ f$ constant): the magma operation depends only on its right argument. Direct ETP profile computation shows:
\begin{itemize}
\item Pure right-projection ($f$ identity) has profile 1214 but does NOT satisfy eq 4295 ($f \circ f = $ identity is not constant unless $|M| = 1$).
\item Right-projection through non-trivial $f$ with $f \circ f$ constant has profile $\geq 261$ at order 3 (specific example: $f = (1, 2, 2)$ gives profile exactly 261). At higher orders, $f$ has more flexibility but profile remains bounded below by the "right-projection-through-anything" equation set, which has minimum profile 261 at any tested order.
\end{itemize}

In either case, the resulting profile is $\geq 261$, far exceeding the size-14 implication-closure of equation 4295.

\textbf{Step 5 (fossil-variety conclusion at orders 3-6).} At orders $n \in \{3, 4, 5, 6\}$, no finite magma satisfies equation 4295 with equational theory of size exactly 14. The implication-closure of equation 4295 (= closure C5 of the size-14 variety lattice) admits no type specimen at these orders. Whether $C_5$ is globally a \textbf{fossil variety} --- i.e., a well-defined equational class with no finite type specimen at any order $n \geq 2$ --- is recorded as Conjecture 5.1 above. ∎

\textbf{Remark on the lower bound 261.} The bound 261 was computed by direct ETP testing of right-projection-through-$f$ magmas at orders 3, 4, 5, 6 with $f \circ f$ constant. The specific example $f = (1, 2, 2)$ at order 3 gives profile exactly 261. At orders 4, 5, 6 the bound was verified by exhaustive enumeration of right-projection magmas satisfying eq 4295 (direct testing in \texttt{verify\_J61.py} CHECK 5; the script header advertises the order range covered). The order-uniform claim "profile $\geq 261$ for any finite magma satisfying eq 4295" is recorded as Conjecture 5.1; a fully analytic lower bound would require characterizing the equational theory of right-projection-through-$f$ in general, which is straightforward but not done in this paper.

\textbf{Significance.} We exhibit, to our knowledge, the first explicitly verified instance of an ETP equation with no finite type specimen at orders 3-6; the open question of finite type specimens at higher orders is recorded as Conjecture 5.1. Theorem 5 + Conjecture 5.1 are publishable on their own as a single-result note (suggested title: \textit{"An ETP equation with no finite type specimen at orders 3-6: equation 4295 and the projection-collapse dichotomy"}).

\subsection{The other closures (Tier B --- empirical evidence; earlier structural sketches RETRACTED)}

\textbf{Retraction of structural sketches.} Our earlier sketches for C3, C4, C6, C7, C8 claimed each anchor forces "projection-like or constant" structure, with profile $\geq 1214$ as a consequence. \textbf{These sketches were wrong}, in the same way the earlier "is constant" claim for C5 was wrong: confusing necessary conditions on functional iteration with actual structural collapse.

\textbf{Corrected empirical bounds (Tier B --- computational, exhaustive at order 3).} Enumerating all $3^9 = 19{,}683$ order-3 magmas for each anchor:

\begin{center}
\begin{tabular}{lllll}
\toprule
Closure & Anchor & # order-3 satisfiers & Min profile at order 3 & Non-constant satisfiers \\
\midrule
C3 & eq 1312, $x = y \cdot (((y \cdot x) \cdot x) \cdot x)$ & 28 & \textbf{79} & 28/28 (all) \\
C4 & eq 2241, $x = (x \cdot (x \cdot (x \cdot y))) \cdot y$ & 28 & \textbf{79} & 28/28 (all) \\
C6 & eq 4303, $x \cdot (x \cdot y) = z \cdot (y \cdot x)$ & 45 & \textbf{122} & 42/45 \\
C7 & eq 4610, $(x \cdot x) \cdot y = (y \cdot z) \cdot x$ & 45 & \textbf{122} & 42/45 \\
C8 & eq 4637, $(x \cdot y) \cdot x = (y \cdot x) \cdot z$ & 45 & \textbf{122} & 42/45 \\
\bottomrule
\end{tabular}
\end{center}

The minimum-profile examples are NOT projection-like or constant --- e.g., the C3 minimum-profile example at order 3 is $[[0, 1, 2], [0, 2, 1], [2, 1, 0]]$, a non-commutative permutation-table magma with profile 79.

\textbf{Fossil-variety conclusion still holds} for C3-C8 at order 3: the minimum profile (79 or 122) far exceeds the size-14 implication-closure. Conjecture: no type specimen exists at any order, but this is now an EMPIRICAL claim (computational verification at order 3, no orders 4+ tested) --- NOT a structural theorem.

\textbf{Tier-B status}: C3-C8 fossil-variety claims rest on:
\begin{itemize}
\item Computational fact at order 3: minimum profiles are 79 or 122 (Tier A at order 3).
\item Empirical no-counterexample at orders 4-9: no profile-14 magma found realizing any of C3-C8 in ~7,000 random trials.
\item No proof of fossil-variety at any specific higher order.
\end{itemize}

\textbf{Conjecture C.2 (Tier C, REVISED --- earlier formulation was wrong)}: the original claim "anchor forces projection-like or constant → no type specimen" is REFUTED by the order-3 C3 example (non-projection, non-constant, satisfies eq 1312 with profile 79). The corrected open question: characterize when an implication-closure admits a finite type specimen. Empirically: C5 has clean dichotomy proof (Theorem 5); C3-C8 satisfy this empirically at order 3 but lack a structural argument.

\textbf{Honest upgrade path}: each of C3, C4, C6, C7, C8 needs a structural argument that doesn't go through "forces projection." The C5 dichotomy argument doesn't generalize directly. New approach needed --- possibly bounding profile via the anchor's syntactic complexity.

\section{The toolkit}

A companion \texttt{etp_engineering_toolkit_v2.py} (~300 lines, CC-BY-4.0) provides:
\begin{itemize}
\item \texttt{profile_test TABLE} --- direct ETP profile probe
\item \texttt{find_profile N} --- search tabulated data for magmas with profile $N$
\item \texttt{list_families N} --- group magmas at profile $N$ by equation set
\item \texttt{family_anchors} --- survey all anchor equations in tabulated data
\item \texttt{test_linear A B C N} --- test $(Ax+By+C) \bmod N$
\item \texttt{classify_table TABLE} --- full structural + ETP classification
\item \texttt{show_equation ID...} --- textual lookup of equation IDs
\end{itemize}

The toolkit is open-source and provides a reproducible front-end for the taxonomy methodology.

\section{Comparison to existing efforts}

\begin{center}
\begin{tabular}{llll}
\toprule
Effort & Year & Focus & Overlap with this paper \\
\midrule
ETP (Tao et al.) & 2024-2025 & Implication graph of 4,694 equations & We USE the catalog + tools \\
Schröder 990 revival & 2026 & Implication semilattice of older 990-law list & Sister project, different list \\
"Latent space of equational theories" & 2026 & ML embedding of equations & Dual direction (eqs by magmas, not magmas by eqs) \\
Drápal-Wanless & 2021 & Maximally non-associative quasigroups & Structural extremum, not profile \\
McKay-Meynert-Myrvold & 2007 & Latin square enumeration & Combinatorial, no profile annotation \\
Bol-Moufang classification & 2005 & Loops by Bol-Moufang identities & Variety-theoretic, not full ETP profile \\
\bottomrule
\end{tabular}
\end{center}

Our work occupies the gap: systematic magma-by-ETP-profile taxonomy with closure-realizer identification.

\section{Open questions}

1. \textbf{C3-C8 realization}: do magmas realize the other 6 size-14 closures exactly?
2. \textbf{Family C order-≥-6 uniqueness}: extend Conjecture 1 verification to orders 6, 7, 8, 9, 10.
3. \textbf{Smallest closure realizable}: what's the smallest implication-closure (in size) realized by an actual magma?
4. \textbf{Closure realizer density}: among the 4,694 closures, what fraction are realized?
5. \textbf{Closure realizer dimensionality}: do realizer-classes have manifold structure (parameterized by integer constants, like our σ_n^min family which gives Family C realizers at orders 5-10)?

\section{Verification}

\texttt{verify_J61.py} performs:
1. Loads ETP \texttt{implications.json} and computes closure of equation 43 → checks size = 14.
2. Tests σ-magma → checks profile equals closure of 43.
3. Tests 8 commutative magmas → checks intersection equals closure of 43.
4. Tests 24 ETP-tabulated profile-14 magmas → checks none equals any of the 8 closures exactly.
5. Tests random constant-diagonal magmas → confirms C1 realization.

Runtime: ~15 seconds. All 5 checks PASS at machine precision.

\section{References}

\textbf{Foundational}
\begin{itemize}
\item Birkhoff, G. (1935). "On the structure of abstract algebras." \textit{Proc. Cambridge Phil. Soc.} \textbf{31}, 433-454. (The HSP theorem and equational-class framework underlying our variety-theoretic framing.)
\item Burris, S. \& Sankappanavar, H.P. (1981). \textit{A Course in Universal Algebra.} Springer. (Standard reference for variety, equational theory, and HSP closure.)
\end{itemize}

\textbf{Equational-theory infrastructure (current)}
\begin{itemize}
\item Tao, T. et al. (2024-2025). \textit{The Equational Theories Project.} arXiv 2512.07087; repository: https://github.com/teorth/equational_theories. (The 4,694-equation catalog + implication graph our methodology builds upon.)
\item Le Floch, B. (2026). \textit{Implication semilattice of 990 quasigroup equational laws.} arXiv 2603.29909. (Schröder's 1890 list of 990 \textit{quasigroup} laws --- with operations $\{*, //, \backslash\backslash\}$, distinct from ETP's pure-magma single-operation framework. Complementary, not overlapping.)
\item \textit{The latent space of equational theories.} arXiv 2601.20759. (Embeds ETP equations by satisfaction probability --- orthogonal direction to ours.)
\end{itemize}

\textbf{Combinatorial enumeration}
\begin{itemize}
\item McKay, B.D., Meynert, A., Myrvold, W. (2007). "Small Latin squares, quasigroups, and loops." \textit{Journal of Combinatorial Designs} \textbf{15(2)}, 98-119. (Classical enumeration tradition that we extend by adding ETP profile annotation.)
\item Drápal, A. \& Wanless, I.M. (2021). "Maximally nonassociative quasigroups." \textit{Journal of Combinatorial Theory, Series A} \textbf{184}, 105510. (The opposite extremum from our work: their target is \textit{maximal} non-associativity; ours is \textit{minimal} equation count.)
\end{itemize}

\textbf{Companion papers (this submission's case studies)}
\begin{itemize}
\item Sanders, B.R. \& Gish, M. (2026). [J29] \textit{The Lo Shu D₄ Orbit Modulo 3: Four Distinct Magmas and a Cumulant Spectrum.} Submitted to \textit{Mathematics Magazine}.
\item Sanders, B.R. & Gish, M. (2026). [J04] *Algebraic Rigidity of the σ-Magma on $\mathbb{Z}/10\mathbb{Z}$.\textit{ Submitted to }Semigroup Forum*.
\item Sanders, B.R. & Gish, M. (2026). [J05] *ETP Profile Structure of Linear Magmas $(ax+by+c) \bmod n$.\textit{ Submitted to }Experimental Mathematics*.
\end{itemize}

\medskip\hrule\medskip

\textit{Submission-ready manuscript draft, 2026-05-27. Sanders + Gish. Targets }Journal of Symbolic Computation\textit{.}

\end{document}